\section{Implementation}
\label{sec:implementation}

\subsection{Technology Stack}

Our platform is implemented using the following technologies:

\begin{table}[htbp]
    \centering
    \caption{Technology stack.}
    \label{tab:tech_stack}
    \begin{tabular}{@{}ll@{}}
        \toprule
        \textbf{Layer} & \textbf{Technology} \\
        \midrule
        Backend API & FastAPI (Python 3.11) \\
        K8s Client & kubernetes-python-client \\
        Database & SQLite + SQLModel \\
        Authentication & JWT (PyJWT) \\
        Frontend & React + TypeScript + Vite \\
        Container Runtime & Docker + containerd \\
        Orchestration & Kubernetes \\
        \bottomrule
    \end{tabular}
\end{table}

\subsection{Backend Components}

The backend consists of several modules:

\begin{itemize}
    \item \textbf{API Layer} (\texttt{api/}): REST API endpoints for container and skill management
    \item \textbf{Service Layer} (\texttt{k8s\_service/}): Kubernetes API integration, Pod exec interface
    \item \textbf{Data Layer} (\texttt{db/}): SQLModel definitions and CRUD operations
\end{itemize}

\subsection{Key Implementation Details}

\subsubsection{Two-Layer Authentication}
The Skill API endpoint enforces two-factor authentication:

\begin{lstlisting}[language=Python]
# Layer 1: Admin API Key
if x_admin_api_key != Config.HAI_K8S_ADMIN_API_KEY:
    raise HTTPException(status_code=403)

# Layer 2: User JWT (extracts user_id, verifies ownership)
user = decode_jwt(authorization)
if container.user_id != user.id:
    raise HTTPException(status_code=403)
\end{lstlisting}

This ensures that only authorized Orchestrators can invoke Skills, and each user can only operate on their own containers.

\subsubsection{Command Execution in Containers}
We use Kubernetes' \texttt{Exec-API} to run commands inside agent containers:

\begin{lstlisting}[language=Python]
from kubernetes.stream import stream as k8s_stream

resp = k8s_stream(
    v1.connect_get_namespaced_pod_exec,
    pod_name, namespace,
    command=["bash", "-c", command],
    stderr=True, stdout=True, stdin=False,
    _preload_content=True, timeout=timeout
)
\end{lstlisting}

The response includes stdout, stderr, and exit code, enabling the Orchestrator to determine success or failure.

\subsubsection{OpenClaw Initialization Pipeline}
The \texttt{openclaw-manager} Skill executes the following steps:

\begin{enumerate}
    \item \textbf{Onboard}: \texttt{openclaw onboard --non-interactive --accept-risk --flow quickstart \ldots}
    \item \textbf{Enable Insecure HTTP}: \texttt{jq '.gateway.controlUi = {"allowInsecureAuth": true}' \ldots}
    \item \textbf{Configure Models}: Fetch template, replace \texttt{\$\{HEPAI\_API\_KEY\}}, write to config
    \item \textbf{Start Gateway}: \texttt{openclaw gateway --port 18789 --bind lan}
\end{enumerate}

\subsection{Deployment}
The platform is deployed on a Kubernetes cluster. Current deployment serves researchers at our institution, with support for GPU-enabled nodes for compute-intensive agent workloads.
